\section{Blind Signature Construction}
\label{sec:blind-signature}

This section presents the vSIS-based blind signature that underlies ARC-SPRUCE. It combines the aligned
linear commitment of Section~\ref{subsec:commitments}, the trapdoor rational-hash interface of
Definition~\ref{def:trapdoor-rational-hash} instantiated with the concrete BB-tran form of
Definition~\ref{def:rational-hash}, and the NIZK-ARK of Definition~\ref{def:nizk-ark}. The signed message is denoted as
\(\vecmu\).

\subsubsection{Randomizing the rational-hash witness}
In the protocol, \(\Holder\) first commits to
\[
(\vecmu,R_0)=(\vecmu,(S\Vert E)),
\]
where \(R_0\) is the linear-side randomness used by the aligned BB-tran target. Concretely, the holder
samples \(S\) and \(E\) from the commitment-randomness distribution and computes
\[
\vecc=\mA_\mu\vecmu+\mA_s S + E.
\]
The issuer then samples admissible inverse-side randomness \(R_1\), and the holder derives the inverse
witness
\[
Z=(B_3-R_1)^{-1}
\]
so that the public target is
\[
T = B_0 + U\vecc + Z = h_{\mathrm{tran},\vecmu,(R_0,R_1)}(B).
\]
The issuing proof therefore binds one hidden opening witness simultaneously to the commitment and to the
public target.

\subsubsection{Blind Signature protocol}
\label{subsec:bs-protocol}
We assume \(\nizkscheme=(\ZKProve,\ZKVerify)\) is a NIZK for the relations below,
\(\ComScheme=(\ComSetup,\Commit,\ComVerify)\) is the aligned linear commitment, and
\[
  \TDRH=(\Setup,\TrapGen,\Eval,\SampPre)
\]
is a trapdoor rational hash according to Definition~\ref{def:trapdoor-rational-hash}. The public setup
samples the aligned commitment key together with a public map \(U\) and public elements \(B_0,B_3\); writing
\[
B_1:=U\mA_\mu,
\qquad
B_2(S\Vert E):=(U\mA_s)S + UE,
\]
this defines the concrete BB-tran target used in issuance and verification.

We define the relations
\[
\relIssueBS:=\left\{ \left(
\begin{pmatrix}
\ComPublicParam,\vecc,T,R_1,U,B
\end{pmatrix},
\begin{bmatrix}
\vecmu\\ S\\ E\\ Z
\end{bmatrix}\right)
~\bigg |~
\begin{matrix}
\ComPublicParam=(\mA_\mu,\mA_s,\BoundMsg) \;\land\\
\vecc = \mA_\mu\vecmu + \mA_s S + E \;\land\\
\vecmu \in \mathcal M \;\land\\
||\vecmu,S,E||_{\infty} \leq \BoundMsg \;\land\\
R_1\in\mathrm{Supp}(\mathcal D_1^{\mathrm{adm}}) \;\land\\
(B_3-R_1)\Had Z = 1 \;\land\\
T = B_0 + U\vecc + Z
\end{matrix}
\right\}.
\]
and
\[
\relSign :=\left\{ \left(
\begin{pmatrix}
\ComPublicParam,\vecmu,\mA,U,B,\BoundMsg,\BoundSig
\end{pmatrix},
\begin{bmatrix}
S\\ E\\ R_1\\ Z\\ \vecu
\end{bmatrix}\right)
~\bigg |~
\begin{matrix}
\ComPublicParam=(\mA_\mu,\mA_s,\BoundMsg) \;\land\\
\vecmu \in \mathcal M \;\land\\
||\vecmu,S,E||_{\infty} \leq \BoundMsg \;\land\\
||\vecu||_\infty \leq \BoundSig \;\land\\
R_1\in\mathrm{Supp}(\mathcal D_1^{\mathrm{adm}}) \;\land\\
(B_3-R_1)\Had Z = 1 \;\land\\
\mA \vecu = B_0 + U(\mA_\mu\vecmu + \mA_s S + E) + Z
\end{matrix}
\right\}.
\]

Our Blind Signature scheme is then defined as follows:
\begin{description}
\item[$\Setup(\SecParam)$:] ~
\begin{enumerate}
\item Run
\[
(\mA_\mu,\mA_s,\BoundMsg)\gets \ComScheme.\ComSetup(\SecParam,\PublicParamsGen).
\]
\item Sample
\[
U\getsunif \Rq^{1\times \ell_{\ComScheme}},\qquad B_0,B_3\getsunif \Rq,
\]
and define the BB-tran description \(B=(B_0,B_1,B_2,B_3)\) by
\[
B_1:=U\mA_\mu,
\qquad
B_2(S\Vert E):=(U\mA_s)S + UE.
\]
\item Set
\[
\ComPublicParam:=(\mA_\mu,\mA_s,\BoundMsg),
\qquad
\mathcal M:=\{\vecmu\in\Rq^{\ell_\mu}: ||\vecmu||_\infty\le \BoundMsg\},
\]
\[
\mathcal X:=\mathcal D_0\times \mathcal D_1^{\mathrm{adm}},
\qquad
\delta:=\delta_{\mathrm{adm}}.
\]
\[
\PublicParamsTrapdoor:=(B,\mathcal M,\mathcal X,\delta,\BoundSig).
\]
\item Output
\[
\PublicParamsBS=(\ComPublicParam,U,\PublicParamsTrapdoor,\BoundMsg).
\]
\end{enumerate}
\item[$\IKeyGen()$:] ~
\begin{enumerate}
\item Compute \((\mA,\trapdoor)\gets \TDRH.\TrapGen(\PublicParamsTrapdoor)\).
\item Output \(\vk=(\mA)\) and \(\sk=(\trapdoor)\).
\end{enumerate}
\item[$\Issue$:] \(\Issuer\) holds \(\sk\) while \(\Holder\) has \(\vecmu,\vk\).
\begin{enumerate}
\item \(\Holder\) samples \(S\gets \chi_s^{\ell_s}\) and \(E\gets \chi_e^{\ell_{\ComScheme}}\).
\item \(\Holder\) computes \(\vecc\gets \mA_\mu\vecmu + \mA_s S + E\) and sends \(\vecc\) to \(\Issuer\).
\item \(\Issuer\) samples \(R_1\gets \mathcal D_1^{\mathrm{adm}}\) and sends \(R_1\) to \(\Holder\).
\item \(\Holder\) sets \(Z\gets (B_3-R_1)^{-1}\) and
\begin{equation}\label{bs:step:issuancenizk}
T\gets B_0 + U\vecc + Z = h_{\mathrm{tran},\vecmu,((S\Vert E),R_1)}(B),
\end{equation}
and produces the issuing NIZK \(\nizkIssue\) using \(\nizkscheme.\ZKProve\) proving knowledge of
\((\vecmu,S,E,Z)\) satisfying relation \(\relIssueBS\).
\item \(\Holder\) sends \((T,\nizkIssue)\) to \(\Issuer\).
\item \(\Issuer\) verifies \(\nizkIssue\) for relation \(\relIssueBS\) on
\((\ComPublicParam,\vecc,T,R_1,U,B)\) using \(\nizkscheme.\ZKVerify\); if it accepts, sample
\(\vecu \gets \TDRH.\SampPre(\trapdoor,T)\) and send \(\vecu\) to \(\Holder\). Otherwise abort.
\item \(\Holder\) checks \(\mA \vecu = T\) and that \( || \vecu ||_\infty \leq \BoundSig \). If not, abort.
\item Set \(\sig:=(\vecu,S,E,R_1)\).
\end{enumerate}
\item[$\Verify(\vk,\vecmu,\sig)$:] For \(\sig=(\vecu,S,E,R_1)\):
\begin{enumerate}
\item Compute \(\vecc\gets \mA_\mu\vecmu + \mA_s S + E\).
\item Reject if \(||\vecmu,S,E||_\infty > \BoundMsg\) or \(||\vecu||_\infty > \BoundSig\).
\item Reject if \(B_3-R_1\notin \Rq^\times\).
\item Compute \(Z\gets (B_3-R_1)^{-1}\) and accept iff
\[
\mA\vecu = B_0 + U\vecc + Z.
\]
\end{enumerate}
\end{description}

\subsubsection{Security}
\begin{proposition}[Correctness]
Assume correctness of \(\nizkscheme,\TDRH\). Then the construction is correct according to
Definition~\ref{def:bs-correctness}.
\end{proposition}

\begin{proof}
The holder computes \(T=B_0+U\vecc+Z\) with \((B_3-R_1)\Had Z=1\), so the issuing relation is satisfiable by
construction. Completeness of the pre-sign NIZK implies that \(\Issuer\) accepts the proof \(\nizkIssue\) for relation
\(\relIssueBS\). By correctness of \(\SampPre\), the returned \(\vecu\) satisfies \(\mA\vecu=T\) and has norm
\(\leq \BoundSig\) except with negligible probability. The output \(\sig=(\vecu,S,E,R_1)\) therefore verifies except
with negligible probability.
\end{proof}

\begin{proposition}[Issuer-view message hiding of issuance]\label{prop:bs-issuance-hiding}
Assume Theorem~\ref{thm:tran-lwe-hiding} and that \(\nizkscheme\) is zero knowledge. Then the public issuance
surface
\[
(\vecc,R_1,T,\nizkIssue)
\]
is computationally hiding with respect to the holder message \(\vecmu\).
\end{proposition}

\begin{proof}
We use two hybrids. In the first hybrid, replace the real issuing proof \(\nizkIssue\) by the simulated proof
generated by the zero-knowledge simulator of \(\nizkscheme\). This changes the issuer view only negligibly. In the
second hybrid, conditioned on the same public parameters and on the public issuer sample \(R_1\), the remaining
holder-to-issuer view is exactly the pair \((\vecc,T)\) from the aligned BB-tran hiding experiment of
Section~\ref{subsec:aligned-hiding}, together with the independent public value \(R_1\). Theorem~\ref{thm:tran-lwe-hiding}
therefore implies that \((\vecc,T)\) is computationally hiding as a function of \(\vecmu\), and adding the independent
public value \(R_1\) preserves that hiding. Combining the two hybrids proves the claim.
\end{proof}

\begin{remark}[Security status of the blind protocol]\label{rem:blind-security-roadmap}
The non-blind hash-and-preimage signature layer underlying this protocol is justified by the DKLW25 reduction
summarized in Lemma~\ref{lem:vsis-unforgeability}. The issuance surface above additionally satisfies the issuer-view
hiding property of Proposition~\ref{prop:bs-issuance-hiding}. A proof of one-more unforgeability would still have to
combine extractability of the pre-sign proof with an interactive hardness statement that genuinely matches the issuance
oracle. Likewise, a full blindness theorem in the sense of Definition~\ref{def:bs-blindness} would still have to
strengthen the issuer-view simulation argument to the complete two-session experiment with correlated challenge
signatures. Accordingly, the formal claims established here are non-blind signature security for the hash-and-preimage
layer and issuer-view message hiding for the issuance surface; one-more unforgeability and full blindness are left
open.
\end{remark}
